\section*{Foreword}\label{sec:foreword}
\addcontentsline{toc}{chapter}{\nameref{sec:foreword}}

My name is Chris Martin, and I am the founding author of \PM, and I am proud to collaborate with some wonderful people to continuously expand its capabilities.  If you're here, you have enough interest in this project to dive into its documentation, so I really appreciate you.  This is the first edition of the version 3 documentation.  If you find mistakes or errors that need to be corrected in the second edition, please don't hesitate to reach out to me\footnote{Christopher R. Martin Ph.D.\\Associate Professor of Mechanical Engineering\\Penn State University; Altoona College\\\href{mailto:crm28@psu.edu}{crm28@psu.edu}}.

If you've derived a latte's worth of enjoyment or extracted a gallon of gasoline (or a liter of petrol) worth of mileage out of the project, consider supporting us to help keep it going.  Just go to \url{http://pyromat.org/support} to see ways to help.

In the meantime, happy calculating!

\section*{Nomenclature}\label{sec:nom}
\addcontentsline{toc}{chapter}{\nameref{sec:nom}}

These variables and units are restricted to derivations in this document.  Because \PM\ has a user-configurable unit system, any of these may be expressed in alternate units.  

The following are parameters with units.

\vspace{1em}
\begin{tabular}{|ccl|}
\hline
Symb. & Units & Description\\
\hline
$a$ & m/s & Speed of sound\\
$c$ & J/kg/K & Specific heat\\
$e$ & J/kg & Internal energy\\
$f$ & J/kg & Free energy\\
$g$ & J/kg & Gibbs energy\\
$h$ & J/kg & Enthalpy\\
$k$ & J/K & Boltzmann constant\\
$M$ & kg & Mass\\
$N$ & count & Number of moles\\
$n$ & m$^{-3}$ & Number density\\
$p$ & Pa & Pressure\\
$q$ & J kg$^{-1}$ & Normalized heat addition\\
$R$ & J kg$^{-1}$ K$^{-1}$ & Gas constant\\
$R_u$ & J kmol$^{-1}$ K$^{-1}$ & Universal gas const.\\
$s$ & J/kg/K & Entropy\\
$T$ & K & Temperature\\
$v$ & m$^3$/kg & Specific volume\\
$W$ & kg/kmol & Molar mass\\
$\rho$ & kg/m$^3$ & Density\\
\hline
\end{tabular}
\vspace{1em}

The following are dimensionless parameters.

\vspace{1em}
\begin{tabular}{|ccl|}
\hline
Symb. & Def. & Description\\
\hline
$\alpha$ & $f / RT$ & ND Free energy\\
$\delta$ & $\rho/\rho_c$ & ND Density\\
$\tau$ & $T_c / T$ & ND Temperature\\
$\gamma$ & $c_p / c_v$ & Sp. heat ratio\\
$y_i$ & $m_i / m$ & Mass fraction\\
$\chi_i$ & $N_i / N$ & Mole fraction\\
$\pi$ & $p / p_c$ & ND Pressure\\
\hline
\end{tabular}
\vspace{1em}

The following are subscripts used throughout the document.

\vspace{1em}
\begin{tabular}{|cl|}
\hline
Symb. & Description\\
\hline
$x_c$ & Critical point\\
$x_f$ & Formation property (e.g. enthalpy of formation)\\
$x_p$ & Constant-pressure (sp.heat)\\
$x_t$ & Triple point\\
$x_v$ & Constant-volume (sp.heat)\\
$x^\circ$ & Property at reference pressure\\
$\overline{x}$ & Property in molar or number units\\
\hline
\end{tabular}

Special constants used in more fundamental calcuations are defined throughout this document.  Many are listed explicitly in Table \ref{tab:constants} in Chapter \ref{ch:units}.
